% Moved from sections/15_des.tex (end-sem split: questions live in the Question Bank PDF).
\section{Data Encryption Standard (DES)}\label{q:15_des}
\noindent{\small Notes: \S\ref{sec:des}.}

\subsection*{Check your understanding}
\begin{enumerate}[leftmargin=1.5em,itemsep=2pt]
  \item \textbf{Why do IP and FP add no security?} They are fixed, keyless and public, so an attacker can apply them himself and attack the 16 rounds directly.
  \item \textbf{Why does round 16 have no swapper?} So that the sequence of operations of the reverse cipher is identical to the cipher; then decryption is the same algorithm with keys $K_{16}, \dots, K_1$.
  \item \textbf{What is the output of S-box 1 for input \ttx{111111}?} Row $(11)_2 = 3$, column $(1111)_2 = 15$; S1[3][15] $= 13 = $ \ttx{1101}.
  \item \textbf{Why does DES need an expansion P-box?} To match the 48-bit round key and to give eight 6-bit S-box inputs; the duplicated edge bits also feed two neighbouring S-boxes, which speeds up the avalanche.
  \item \textbf{How many bits are shifted in total, and why does it matter?} $4 \times 1 + 12 \times 2 = 28$, a full rotation, so $C_{16} = C_0$ and $D_{16} = D_0$ and decryption keys can be generated by right shifts from the original halves.
  \item \textbf{Is the DES S-box invertible? Does it matter?} No, a $6 \to 4$ box has 4 preimages per output; it does not matter because it sits inside $f$, which a Feistel cipher never inverts.
  \item \textbf{What is special about a weak key?} All 16 round keys are equal, so $E_K(E_K(P)) = P$; there are 4 weak keys, 6 pairs of semi-weak keys and 48 possible weak keys.
  \item \textbf{Why is double DES not 112-bit secure?} Meet-in-the-middle: with one known $(P, C)$, tabulate $E_{K_1}(P)$ for all $2^{56}$ keys and match against $D_{K_2}(C)$ for all $2^{56}$ keys, which costs $2^{57}$ operations instead of $2^{112}$.
  \item \textbf{Why is triple DES EDE and not EEE?} With $K_1 = K_2$ (or $K_2 = K_3$), EDE collapses to single DES, so new equipment stays compatible with single-DES equipment.
  \item \textbf{State the complement property and its consequence.} $E_{\bar K}(\bar P) = \overline{E_K(P)}$, so a chosen-plaintext brute-force attack tests two keys per encryption and needs only $2^{55}$ trials.
\end{enumerate}
